\documentclass[preprint,12pt]{elsarticle}
\biboptions{numbers,sort&compress}

\usepackage{amsmath,amssymb,amsfonts}
\usepackage{graphicx}
\usepackage{booktabs}
\usepackage{multirow}
\usepackage{url}
\usepackage{hyperref}
\usepackage{algorithm}
\usepackage{algorithmic}
\usepackage{subcaption}
\usepackage{microtype}
\usepackage{tabularx}
\usepackage{lineno}
\usepackage{tikz}
\usetikzlibrary{shapes.geometric, arrows.meta, positioning, calc, shadows, fit, backgrounds}

% Custom Palette for Publication-Grade TikZ Architecture
\definecolor{bannerbg}{RGB}{24, 33, 47}
\definecolor{cblue}{RGB}{37, 99, 235}
\definecolor{cbluebg}{RGB}{239, 246, 255}
\definecolor{cteal}{RGB}{13, 148, 136}
\definecolor{ctealbg}{RGB}{240, 253, 250}
\definecolor{cpurple}{RGB}{124, 58, 237}
\definecolor{cpurplebg}{RGB}{245, 243, 255}
\definecolor{cgreen}{RGB}{22, 163, 74}
\definecolor{cgreenbg}{RGB}{240, 253, 244}
\definecolor{camber}{RGB}{217, 119, 6}
\definecolor{camberbg}{RGB}{255, 251, 235}
\definecolor{cred}{RGB}{220, 38, 38}
\definecolor{credbg}{RGB}{254, 242, 242}
\definecolor{ctxtdark}{RGB}{15, 23, 42}
\definecolor{ctxtgray}{RGB}{71, 85, 105}

\journal{Sustainable Computing: Informatics and Systems}

\begin{document}

\begin{frontmatter}

\title{SHIFT-GPU: Shift-Aware, Risk-Calibrated and Physics-Guided Learning for Reliable LLM Inference Hardware Selection}

\author[inst1]{Jason Pandian\corref{cor1}}
\ead{pandiajason@gmail.com}

\cortext[cor1]{Corresponding author. Department of Information Technology, Nehru Institute of Technology, Coimbatore, Tamil Nadu, 641105, India.}

\affiliation[inst1]{organization={Department of Information Technology, Nehru Institute of Technology},
            addressline={Thirumalayampalayam}, 
            city={Coimbatore},
            postcode={641105}, 
            state={Tamil Nadu}, 
            country={India}}

\begin{abstract}
Selecting optimal hardware accelerators for Large Language Model (LLM) serving in heterogeneous cloud clusters is a critical operational problem in sustainable computing. In production data centers, exhaustive empirical profiling across models, batch sizes, and GPU generations is environmentally and financially prohibitive. While metadata-driven predictors estimate power draw and latency, existing models suffer severe performance degradation when evaluating unseen GPU architectures, causing Service Level Agreement (SLA) violations and excessive energy consumption. To resolve this, we present \textbf{SHIFT-GPU}, a framework uniting physics-guided deep residual ensembles, heteroscedastic and epistemic covariate-shift calibration, and constrained admission control. At its core is \textbf{PICAN} (\textit{Physics-Informed Cross-Attention Network}), which couples a deterministic microarchitectural Roofline latency bound with decoupled workload--hardware latent encoders ($w_x, h_g \in \mathbb{R}^d$), smooth multiplicative exponential residuals ($\hat{T} = T_{\text{roof}} \cdot \exp(\mathrm{softplus}(r_\theta))$), and Gaussian Negative Log-Likelihood uncertainty estimation across a 5-seed ensemble. To manage microarchitectural domain shift, SHIFT-GPU incorporates shift-adjusted conformal upper bounds scaled by Mahalanobis hardware discrepancy. We also evaluate few-shot FiLM calibration ($k \in \{1, 3, 5\}$) across probing heuristics, finding that uniform random probing ($7.37\%$--$7.38\%$ MdAPE) outperforms boundary-targeted heuristics ($7.93\%$--$8.10\%$) while zero-shot 5-seed ensembles already anchor global generalization. Evaluated under a NaN-blinded zero-leakage protocol across 4,163 production configurations spanning 42 open-source LLMs (0.12B--27B parameters) and 8 NVIDIA GPU architectures, SHIFT-GPU reduces Leave-One-GPU-Out (LOGO) inter-token latency error from 11.98\% (matched XGBoost rerun) and 10.68\% (Roofline-augmented XGBoost) down to \textbf{6.83\% MdAPE} (relative reductions of 43.0\% and 36.1\%, respectively; single-seed mean 7.22\% $\pm$ 0.75\%), achieves 3.31\%--6.55\% error across modern accelerators (L4, H100, L40S, A30, H200), and reaches \textbf{6.21\% MdAPE} under 42-model Leave-One-LLM-Out (LOLO) generalization (vs.\ 7.72\% for XGBoost). Under production cluster selection across all evaluated SLAs ($T_{\max} \in \{20, 25, 30, 35, 40\}$~ms), risk-calibrated SHIFT-GPU records \textbf{0 violations out of 125 accepted workloads (0.0\%, 0 abstentions)}, whereas naive Minimum TDP incurs 31.2\%--51.2\% violations, Roofline-only Minimum TDP incurs 5.6\%--24.8\% violations, and uncalibrated point predictors incur up to 2.4\% violations.
\end{abstract}

%% Research highlights
\begin{highlights}
\item Cuts unseen GPU latency error from 11.98\% to 6.83\% (-43.0\%) over XGBoost.
\item Outperforms Roofline-augmented XGBoost (10.68\% to 6.83\%, -36.1\%) across 8 GPUs.
\item Achieves 6.21\% LOLO error across 42 LLMs under strict NaN-blinded verification.
\item Conformal bounds maintain 0/125 SLA violations across 20--40 ms latency thresholds.
\item Reproducible open-source pipeline, 5-seed benchmarks, and tables publicly released.
\end{highlights}

\begin{keyword}
Large Language Models \sep GPU Infrastructure \sep Physics-Informed Neural Networks \sep Conformal Prediction \sep Active Learning \sep Sustainable Computing.
\end{keyword}

\end{frontmatter}

\linenumbers

\section{Introduction}
\label{sec:intro}

The exponential growth of generative artificial intelligence and Large Language Models (LLMs) has introduced unprecedented computational and environmental demands on cloud data center infrastructure \citep{strubell2019energy, patterson2021carbon, suscom_desislavov2023}. In contemporary enterprise deployments, inference represents the dominant fraction of the total operational expenditure and carbon footprint over a model's operational lifecycle. Serving autoregressive transformer architectures requires maintaining strict Service Level Agreements (SLAs)---typically specified as an upper bound on mean Inter-Token Latency (ITL) or Time to First Token (TTFT)---while simultaneously minimizing total cluster energy consumption and operational power draw \citep{kwon2023vllm, suscom_malekloo2018}.

To optimize throughput per Watt, modern cloud providers deploy heterogeneous GPU clusters spanning multiple architectural generations, including NVIDIA Hopper (H100, H200), Ada Lovelace (L40S, L4), Ampere (A100, A30), Turing (T4), and Volta (V100). When an infrastructure operator seeks to deploy a new LLM workload, they must decide: \textit{Which GPU candidate satisfies the required latency SLA with the lowest operational power draw?}

In the sustainable computing literature, optimizing data center hardware allocation and operational energy has evolved from virtualized machine consolidation \citep{suscom_malekloo2018, suscom_kumar2022} and server node-level power capping coordination \citep{suscom_powercoord2020} toward multi-objective optimization for ML inference configurations \citep{suscom_multiobj2026} and predictive perception layers for LLM serving on high-performance computing clusters \citep{suscom_forecasting2026}. Furthermore, queueing-theoretic studies in this journal have revealed that aggressive power minimization risks severe tail latency spikes when operational envelopes are improperly matched to workload demands \citep{suscom_tradeoffs2025}, while macroscopic analyses of AI inference energy have proven that naive parameter scaling fails to capture memory-bandwidth bottlenecks \citep{suscom_desislavov2023}.

Traditionally, answering the hardware selection question requires exhaustive physical profiling: systematically deploying the LLM across every candidate accelerator under varying prompt lengths, concurrency arrival rates, and batch sizes. However, with dozens of active foundation models and hundreds of potential accelerator configurations, empirical profiling incurs massive energy consumption, creates severe scheduling overhead, and delays workload onboarding.

To bypass exhaustive profiling, recent work has explored metadata-based regression models. Notably, WattGPU \citep{fadel2026wattgpu} demonstrated that tree ensembles (XGBoost) trained on public LLM parameters and static GPU specifications can predict inference power and latency. However, as demonstrated in our empirical baseline analysis, existing metadata models suffer from a fundamental failure mode: \textbf{microarchitectural covariate shift on unseen hardware}. When evaluated in-distribution (standard 5-fold cross-validation), tree models achieve an impressive Median Absolute Percentage Error (MdAPE) of 2.84\% for latency and 2.45\% for power. Yet, when evaluated under Leave-One-GPU-Out (LOGO)---simulating the arrival of an unfamiliar hardware architecture---the latency prediction error surges by over 300\% to 11.98\% (and exceeds 17\% on modern architectures such as L4 and A30).

In a production scheduling environment, this degradation is severe. Because point-estimate models lack calibrated uncertainty bounds and microarchitectural inductive biases, they frequently recommend under-provisioned GPUs that suffer uncalibrated SLA violations or over-provisioned accelerators that waste substantial energy.

\subsection{Research Contributions}
To resolve this fundamental gap, this paper presents \textbf{SHIFT-GPU}, an end-to-end computer science framework for shift-aware, risk-calibrated, and physics-guided hardware selection. The core scientific contributions of this work are:

\begin{enumerate}
    \item \textbf{PICAN Architecture (Physics-Informed Cross-Attention Network):} We design a neural architecture that integrates a deterministic microarchitectural Roofline latency lower bound ($T_{\text{roof}} = \max(2 N_{\text{params}}/(B_{\text{mem}} \cdot B(s)),\, 2 N_{\text{params}}/F_{\text{peak}})$) into the hypothesis space. By decoupling the deterministic physical memory-bandwidth streaming floor from learned residual scaling via smooth multiplicative exponential residuals ($\hat{T} = T_{\text{roof}} \cdot \exp(\mathrm{softplus}(r_\theta))$) and Gaussian NLL variance heads across a 5-seed ensemble, PICAN enforces physical consistency across orders of magnitude of latency.
    \item \textbf{Empirical Analysis of Bipartite vs.\ Flat Neural Roofline Encoding:} We evaluate bipartite dual-encoder coupling ($Q = W_q w_x, K = W_k h_g, V = W_v h_g$) alongside both a flat concatenated MLP + Roofline neural baseline and a Roofline-augmented XGBoost baseline ($\log(T / T_{\text{roof}})$ target). We prove mathematically and confirm empirically that because pairwise evaluation uses singleton hardware tokens ($L_k = 1$), softmax degenerates to $1.0$: both PICAN (6.83\% MdAPE, 0.0\% SLA violations) and concatenated MLP+Roofline (6.42\% MdAPE, 0.8\% SLA violations) fall within single-seed variance ($\pm 0.75\%$), whereas both substantially outperform Roofline-augmented XGBoost (10.68\% MdAPE) and raw XGBoost (11.98\% MdAPE). This isolates smooth neural Roofline-residual modeling as the true driver of cross-GPU generalization over axis-aligned trees.
    \item \textbf{Covariate-Shift Conformal Risk Calibration:} We formalize shift-adjusted conformal prediction driven by learned heteroscedastic NLL variance plus deep-ensemble epistemic variance, scaled by microarchitectural Mahalanobis distance in hardware specification space, bounding tail latency risk during scheduling.
    \item \textbf{Transparent Evaluation of Few-Shot FiLM Adaptation:} We benchmark few-shot FiLM parameter calibration ($k \in \{1, 3, 5\}$) without retraining backbone weights across uniform random, K-Means diversity, maximal uncertainty, and SLA-boundary (ROAA) probing strategies. We report the clear empirical finding that when a well-calibrated 5-seed zero-shot ensemble already achieves 6.83\% MdAPE, small-shot gradient updates slightly perturb global MdAPE (7.37\%--8.10\%) while improving linear parity ($r$ from 0.995 to 0.997), with Uniform Random proving most robust.
    \item \textbf{Comprehensive Zero-Leakage LOGO and LOLO Validation:} We conduct an exhaustive evaluation across 4,163 benchmark configurations spanning 42 open-source LLMs (0.12B--27B parameters) and 8 NVIDIA GPU architectures under an automated NaN-blinded zero-leakage protocol where blinded predictions match identically. Results demonstrate that SHIFT-GPU reduces unseen-GPU error from 11.98\% (matched XGBoost) and 10.68\% (XGBoost + Roofline) down to \textbf{6.83\% MdAPE} (43.0\% and 36.1\% relative reductions), achieves \textbf{6.21\% MdAPE} in 42-model Leave-One-LLM-Out (LOLO) generalization, and eliminates SLA violations (\textbf{0 / 125 accepted workloads}) across all latency bounds $T_{\max} \in [20, 40]$~ms.
\end{enumerate}

The remainder of this paper is structured as follows. Section \ref{sec:related} reviews related work in sustainable computing. Section \ref{sec:formulation} presents the formal problem formulation. Section \ref{sec:method} details the SHIFT-GPU methodology, PICAN architecture, and software design. Section \ref{sec:experiments} presents the experimental setup, while Section \ref{sec:results} discusses the empirical results (E1--E7). Section \ref{sec:discussion} analyzes cloud systems implications, Section \ref{sec:threats} addresses threats to validity, and Section \ref{sec:conclusion} concludes.

\section{Related Work}
\label{sec:related}

\subsection{Inference Sustainability and Resource Management in Sustainable Computing}
Sustainable computing literature has extensively studied the dual challenges of performance assurance and energy minimization in cloud infrastructure. Foundational work by Malekloo et al. \citep{suscom_malekloo2018} established the mathematical principles of SLA-compliant, energy-efficient resource allocation, formulating constrained optimization models that balance aggressive power minimization against SLA violation penalties in virtualized environments. Expanding upon these principles, Kumar et al. \citep{suscom_kumar2022} provided a comprehensive taxonomy of machine learning methods for cloud workflow scheduling, virtual machine placement, and energy-aware server consolidation, explicitly highlighting that black-box regression models are vulnerable to severe prediction collapse when operating outside their training distributions. At the physical node level, Azimi et al. \citep{suscom_powercoord2020} introduced PowerCoord, demonstrating coordinated dynamic power capping across heterogeneous multi-CPU and multi-GPU architectures using reinforcement learning under thermal design boundaries.

With the advent of deep learning and Large Language Models, the relationship between hardware utilization, power draw, and execution latency became governed by complex microarchitectural constraints. Desislavov et al. \citep{suscom_desislavov2023} conducted an exhaustive empirical analysis of AI inference energy consumption across diverse model architectures and hardware platforms, demonstrating that classical parameter-versus-performance scaling laws fail to capture non-linear memory and computational saturation. Specifically, they revealed that operational energy efficiency is dominated by arithmetic intensity and memory bandwidth bottlenecks rather than parameter count alone. Addressing the operational selection problem, Casta{\~n}o et al. \citep{suscom_multiobj2026} proposed data-driven multi-objective optimization frameworks to navigate Pareto trade-offs across energy consumption, execution performance, and financial cost during ML inference configuration selection. However, their optimization search was restricted to pre-profiled, static hardware platforms and explicitly identified zero-shot transfer to newly released hardware accelerators as a vital open problem. Complementing these empirical systems, Yao et al. \citep{suscom_tradeoffs2025} investigated power and response-time trade-offs from an analytical queueing perspective, rigorously proving that aggressive power capping creates tail latency instability unless buffer margins are calibrated to execution variability. Recently, Aguirre-Usandizaga et al. \citep{suscom_forecasting2026} developed a forecasting-based perception layer for energy-aware resource management in LLM serving on HPC clusters, establishing that proactive predictive modeling is indispensable for mitigating cluster over-provisioning and idle energy waste.

The SHIFT-GPU framework presented in this paper directly emerges from and unifies this rich body of research published in \textit{Sustainable Computing: Informatics and Systems}. Specifically, SHIFT-GPU resolves the critical open challenges highlighted by this lineage: (1) we translate the memory-bandwidth bottleneck insights of Desislavov et al. \citep{suscom_desislavov2023} into an inductive neural network architecture by embedding an analytical Roofline memory-streaming floor ($T_{\text{roof}}$); (2) we solve the out-of-domain hardware generalization challenge posed by Casta{\~n}o et al. \citep{suscom_multiobj2026} and Aguirre-Usandizaga et al. \citep{suscom_forecasting2026} through microarchitectural bipartite cross-attention; (3) we address the distribution-shift vulnerability identified by Kumar et al. \citep{suscom_kumar2022} using group-conditional conformal calibration; and (4) we operationalize the queueing power-latency trade-offs of Yao et al. \citep{suscom_tradeoffs2025} and the SLA-compliance paradigm of Malekloo et al. \citep{suscom_malekloo2018} by enforcing calibrated tail-latency bounds to eliminate SLA violations during energy-minimizing hardware selection.

\subsection{LLM Serving Systems and Performance Characterization}
Autoregressive LLM inference consists of two distinct operational phases: the compute-bound prefill phase (processing prompt tokens in parallel) and the memory-bandwidth-bound decode phase (generating tokens sequentially) \citep{kwon2023vllm}. Modern serving frameworks employ specialized runtime mechanisms to maximize hardware saturation: vLLM \citep{kwon2023vllm} introduces PagedAttention and continuous batching for efficient KV-cache memory management, while DeepSpeed-Inference \citep{aminabadi2022deepspeed} leverages customized transformer execution kernels, multi-GPU tensor slicing, and memory offloading. Despite these software advances, the underlying hardware efficiency is dictated by the classical Roofline model \citep{williams2009roofline}, which defines the operational boundaries between compute-limited and bandwidth-limited execution. WattGPU \citep{fadel2026wattgpu} demonstrated that tree ensembles (XGBoost) trained on public LLM parameters and static GPU specifications can predict inference power and latency across unseen GPUs. In our matched rerun under Leave-One-GPU-Out (LOGO) cross-validation without hardware/model identifier columns, XGBoost achieves 11.98\% median absolute percentage error for inter-token latency (10.68\% even when trained on Roofline residual targets). Because standard tree models lack smooth microarchitectural interpolation and calibrated uncertainty intervals, their point estimates risk triggering uncalibrated SLA violations during deployment.

\subsection{Conformal Prediction and Covariate Shift}
Conformal prediction provides a distribution-free framework for constructing valid predictive intervals with finite-sample marginal coverage guarantees under exchangeability \citep{angelopoulos2023conformal}. Standard split conformal prediction assumes that calibration and test data are drawn exchangeably from the same distribution ($P_{\text{cal}} = P_{\text{test}}$). When evaluating an unseen accelerator architecture, exchangeability cannot be assumed because the test GPU represents a distinct hardware population. \citet{tibshirani2019conformal_covariate} established the theoretical foundations of weighted conformal prediction under covariate shift using likelihood-ratio density re-weighting. However, likelihood-ratio density estimation cannot be reliably evaluated across sparse, discrete microarchitectural feature spaces. In SHIFT-GPU, we employ an empirical Mahalanobis distance heuristic to scale interval widths under hardware discrepancy. Because exchangeability does not hold across distinct GPU populations and the Mahalanobis adjustment is an empirical heuristic rather than a formal theoretical guarantee, we treat our calibrated bounds as empirical and report observed coverage on held-out hardware.

\subsection{Active Learning and Few-Shot Adaptation}
Active learning aims to maximize predictive accuracy while minimizing the labeling or profiling budget \citep{settles2009active}. While pool-based active learning traditionally queries samples with maximal epistemic variance, such heuristics are sub-optimal for operational decisions with asymmetric loss functions. In SHIFT-GPU, we formulate Regret-Optimal Active Acquisition (ROAA) to query hardware configurations specifically situated near the critical SLA feasibility boundary.

\section{Problem Formulation}
\label{sec:formulation}

\subsection{Inference Workload and Hardware Specification Spaces}
Let an LLM inference workload be defined as a tuple $x = (m, s) \in \mathcal{X}$, where $m \in \mathcal{M}$ represents the model architecture metadata and $s \in \mathcal{S}$ represents the operational serving scenario. Specifically, $m$ includes parameter count $N_{\text{params}}$, hidden dimension $H$, layer count $L$, attention head count $N_h$, key-value head count $N_{kv}$, and KV-cache footprint $C_{kv}$. The operational scenario $s \in \{\text{offline}, \text{server\_low}, \text{server\_high}\}$ specifies the request arrival process and batching dynamics.

Let a candidate hardware accelerator be defined as a vector $g \in \mathcal{G}$, characterized by microarchitectural parameters: memory bandwidth $B_{\text{mem}}$ (GB/s), peak FP16 compute capacity $F_{\text{peak}}$ (TFLOPS), device memory capacity $M_{\text{size}}$ (GB), Thermal Design Power $TDP$ (Watts), base clock frequency $f_{\text{base}}$ (MHz), boost ratio $\kappa_{\text{boost}}$, transistor count $C_{\text{trans}}$, and architectural family generation $A_{\text{gen}} \in \{\text{Volta}, \text{Turing}, \text{Ampere}, \text{Ada}, \text{Hopper}\}$.

For any workload-hardware pair $(x, g)$, the system exhibits two ground-truth execution properties:
\begin{enumerate}
    \item Inter-Token Latency (ITL): $T(x, g) \in \mathbb{R}^+$ (seconds per token).
    \item Mean Operational Power Draw: $P(x, g) \in \mathbb{R}^+$ (Watts).
\end{enumerate}

\subsection{The Constrained Hardware Selection Problem}
Following the SLA-compliant resource allocation principles formalized by Malekloo et al. \citep{suscom_malekloo2018} and the Pareto optimization formulation of Casta{\~n}o et al. \citep{suscom_multiobj2026}, we formulate accelerator selection as a constrained optimization problem. Given a target workload $x \in \mathcal{X}$ and a latency Service Level Agreement threshold $T_{\max} \in \mathbb{R}^+$, the cluster scheduler's objective is to select the accelerator $g^* \in \mathcal{G}$ that minimizes operational power consumption while satisfying the latency constraint with confidence $1 - \alpha$:
\begin{equation}
\label{eq:opt_problem}
g^* = \arg\min_{g \in \mathcal{G}} P(x, g) \quad \text{subject to} \quad \mathbb{P}\left(T(x, g) \le T_{\max}\right) \ge 1 - \alpha
\end{equation}

When ground truth measurements $T(x, g)$ and $P(x, g)$ are unavailable, the selection must rely on surrogate predictors $\hat{T}(x, g)$ and $\hat{P}(x, g)$. If a model relies purely on point estimates $\hat{T}(x, g) \le T_{\max}$, any prediction underestimation ($\hat{T} < T \le T_{\max}$) results in an operational SLA violation. 

\subsection{Operational Regret Formulation}
To rigorously quantify the cost of suboptimal hardware selection, we formalize the operational deployment regret $\mathcal{R}(g^* \mid x)$ relative to the true empirical Oracle accelerator $g_{\text{oracle}}$:
\begin{equation}
\label{eq:regret}
\mathcal{R}(g^* \mid x) = \max\left(0, P(x, g^*) - P(x, g_{\text{oracle}})\right) + \lambda_{\text{SLA}} \cdot P(x, g^*) \cdot \mathbb{I}\left(T(x, g^*) > T_{\max}\right)
\end{equation}
where $\mathbb{I}(\cdot)$ is the indicator function and $\lambda_{\text{SLA}} \ge 1$ is an asymmetric penalty factor ($\lambda_{\text{SLA}} = 5.0$ in our benchmarks) reflecting the operational penalty of an SLA breach. If a risk-calibrated policy abstains from admission despite at least one feasible accelerator existing under oracle knowledge ($\mathcal{G}_{\text{feas}}^{\text{oracle}} \neq \emptyset$), a fixed abstention regret penalty of 50.0 is assessed.

\section{The SHIFT-GPU Methodology}
\label{sec:method}

\begin{figure*}[t]
\centering
\includegraphics[width=\textwidth]{figures/fig1_architecture.pdf}
\caption{Logical abstract architecture and constrained decision pipeline of the SHIFT-GPU framework. (1) Latent encoding maps workload $x$ and hardware candidate $g$ into latent spaces $w_x, h_g \in \mathbb{R}^d$ under zero-leakage blinding; (2) Bipartite interaction couples embeddings via multi-head cross-attention, residual layer normalization, and FiLM latent modulation ($z^* = \gamma_k \odot z + \beta_k$); (3) Physics recombination enforces the analytical Roofline floor ($T_{\text{roof}}$) and thermal envelope ($\mathrm{TDP}$) via multiplicative neural residual heads; and (4) Risk calibration and selection evaluates covariate shift $d_{\text{shift}}(g)$, bounds tail latency via conformal calibration $\hat{T}_{\text{upper}}$, and identifies the power-minimal feasible accelerator $g^* = \arg\min_{g \in \mathcal{G}_{\text{feas}}} \hat{P}(x, g)$ with active probe feedback.}
\label{fig:architecture}
\end{figure*}

The SHIFT-GPU framework comprises five modular, mathematically coupled components designed to solve Eq.~\eqref{eq:opt_problem}, illustrated in Figure \ref{fig:architecture}.

\subsection{Component A: Deterministic Microarchitectural Roofline Base}
Rather than treating latency as an unconstrained regression target, we embed the physical limits of GPU execution into the hypothesis space, directly operationalizing the empirical findings of Desislavov et al. \citep{suscom_desislavov2023} regarding arithmetic intensity and memory-bandwidth saturation. During the decode phase of autoregressive generation, generating each token requires reading model parameters across the memory bus once, while computing $2 N_{\text{params}}$ floating-point operations.

Accounting for sequence parallelism in continuous batching, the effective batch size is denoted as $B(s)$, where $B(\text{server}) = 1.0$ and $B(\text{offline}) = 256.0$ (reflecting the default concurrency saturation cap in modern serving systems such as vLLM). The deterministic microarchitectural Roofline lower bound $T_{\text{roof}}(x, g)$ is defined as:
\begin{equation}
\label{eq:roofline}
T_{\text{roof}}(x, g) = \max\left( \frac{2 \cdot N_{\text{params}}}{B_{\text{mem}}(g) \cdot B(s)}, \frac{2 \cdot N_{\text{params}}}{F_{\text{peak}}(g)} \right)
\end{equation}

For power draw, the physical operational envelope is bounded from above by the manufacturer's Thermal Design Power ($TDP(g)$).

\subsection{Component B: PICAN (Physics-Informed Cross-Attention Network)}
To learn the complex non-linear residual dynamics---including kernel scheduling overheads, KV-cache indexing stalls, and memory channel contention---we design \textbf{PICAN}:

\begin{enumerate}
    \item \textbf{Workload Encoder $E_W$:} Maps continuous workload parameters $x_{\text{cont}}$ and categorical embeddings (scenario, model architecture) to a dense workload latent state $w_x \in \mathbb{R}^d$:
    \begin{equation}
    w_x = \text{MLP}_W\left( [x_{\text{cont}} \parallel \text{Embed}(s) \parallel \text{Embed}(A_{\text{model}})] \right)
    \end{equation}
    \item \textbf{Hardware Encoder $E_H$:} Maps microarchitectural hardware specifications $g_{\text{cont}}$ and categorical embeddings (memory type, architectural family) to a hardware latent state $h_g \in \mathbb{R}^d$:
    \begin{equation}
    h_g = \text{MLP}_H\left( [g_{\text{cont}} \parallel \text{Embed}(M_{\text{type}}) \parallel \text{Embed}(A_{\text{gen}})] \right)
    \end{equation}
    \item \textbf{Microarchitectural Bipartite Cross-Attention (MCA):} Rather than concatenating heterogeneous workload and hardware features into a single unconstrained MLP, PICAN projects workload features and hardware specifications into distinct embedding spaces via dedicated encoders $E_W$ and $E_H$. A bipartite cross-attention layer couples the workload query with hardware keys and values:
    \begin{equation}
    Q = W_q w_x, \quad K = W_k h_g, \quad V = W_v h_g
    \end{equation}
    \begin{equation}
    z_{\text{mca}} = \text{Softmax}\left(\frac{Q K^T}{\sqrt{d}}\right) V, \quad z = \text{LayerNorm}(w_x + z_{\text{mca}})
    \end{equation}

    \textit{Mathematical Qualification on Pairwise Attention Dynamics:} We explicitly qualify the operational mechanics of Eq.~(6)--(7) under pairwise evaluation. In pairwise training and inference, each forward pass evaluates a single workload specification $x$ against a single candidate accelerator $g$ ($L_q = 1, L_k = 1$). For a scalar dot-product with key length unity, the softmax normalization evaluates identically to $\mathrm{Softmax}([s]) = 1.0$ for any real scalar $s = Q K^\top / \sqrt{d}$. Consequently, $z_{\text{mca}} = V = W_v h_g$, and gradients to $W_q$ and $K = W_k h_g$ through the softmax function vanish under single-pair inputs. Under this regime, the MCA block operates as a decoupled projection and residual fusion layer ($z = \text{LayerNorm}(w_x + W_o W_v h_g)$). As confirmed empirically in Ablation B (Section \ref{sec:results}), when both models share the same smooth Roofline objective and 5-seed NLL training, concatenated MLP+Roofline achieves 6.42\% MdAPE (0.8\% SLA violations) compared to dual-encoder PICAN's 6.83\% MdAPE (0.0\% SLA violations, within single-seed variance $\pm 0.75\%$). Thus, the cross-GPU performance gain over XGBoost (11.98\%) and Roofline-augmented XGBoost (10.68\%) stems primarily from smooth neural residual interpolation over the deterministic Roofline floor rather than dynamic attention weighting. Extending PICAN to joint multi-GPU list-wise encoding ($L_k = |\mathcal{G}| > 1$) represents a natural path to unlock true workload-dependent attention weights.
    \item \textbf{Smooth Multiplicative Physics-Residual Head:} Rather than unconstrained additive correction, latency prediction is formulated as an exponential multiplicative scaling over the deterministic Roofline floor:
    \begin{equation}
    \label{eq:mult_residual}
    \hat{T}(x, g) = T_{\text{roof}}(x, g) \cdot \exp\left(\mathrm{softplus}(r_\theta(z))\right)
    \end{equation}
    where $r_\theta(z) \in \mathbb{R}$ represents the learned non-dimensional residual overhead logit. Because $\mathrm{softplus}(r_\theta) = \log(1 + e^{r_\theta}) > 0$ is smooth and strictly positive, Eq.~\eqref{eq:mult_residual} preserves non-vanishing gradients and guarantees $\hat{T}(x, g) > T_{\text{roof}}(x, g)$.
    
    Similarly, for operational power draw, predictions are bounded within the physical TDP envelope:
    \begin{equation}
    \hat{P}(x, g) = TDP(g) \cdot \sigma\left(p_\theta(z)\right)
    \end{equation}
    where $\sigma(\cdot)$ denotes the sigmoid function, guaranteeing $\hat{P} \in (0, TDP(g))$.
    \item \textbf{Heteroscedastic NLL + Deep-Ensemble Predictive Uncertainty:} Alongside point estimates, PICAN outputs log-variance $s_\theta(z) \in [-6, 4]$ trained jointly with Smooth L1 residual regression and Gaussian Negative Log-Likelihood (NLL) across $M=5$ independent random seeds ($\{42, 123, 456, 789, 2026\}$) using cosine learning-rate annealing (150 epochs):
    \begin{equation}
    \mathcal{L}_{\text{ITL}} = \mathrm{SmoothL1}_{\beta=0.05}\!\left(\tilde{r}_\theta, r_{\text{true}}\right) + \frac{1}{10}\,\mathbb{E}\!\left[\frac{1}{2} e^{-s_\theta} (r_{\text{true}} - \tilde{r}_\theta)^2 + \frac{1}{2} s_\theta\right]
    \end{equation}
    where $\tilde{r}_\theta = \mathrm{softplus}(r_\theta)$ and $r_{\text{true}} = \log(T / T_{\text{roof}})$. For ensemble inference across the $M$ seeds, the aggregate latency prediction and total uncertainty combine aleatoric variance and inter-seed epistemic disagreement in linear latency space:
    \begin{align}
    \hat{T}(x, g) &= T_{\text{roof}}(x, g)\cdot \exp\!\left(\frac{1}{M}\sum_{m=1}^M \tilde{r}_\theta^{(m)}\right), \\
    \hat{\sigma}_{\text{total}}(x, g) &= \hat{T}(x, g)\sqrt{\frac{1}{M}\sum_{m=1}^M e^{s_\theta^{(m)}} + \mathrm{Var}_m\!\left(\tilde{r}_\theta^{(m)}\right)}
    \end{align}
\end{enumerate}

\subsection{Component C: Microarchitectural Covariate Shift Quantification}
When an unseen GPU $g_{\text{target}}$ is introduced, its distance from the training distribution manifold $\mathcal{D}_{\text{train}}$ dictates the severity of expected covariate shift. We quantify this discrepancy using the microarchitectural Mahalanobis distance in normalized hardware feature space:
\begin{equation}
\label{eq:mahalanobis}
d_{\text{shift}}(g_{\text{target}}, \mathcal{D}_{\text{train}}) = \sqrt{(g_{\text{target}} - \mu_{\text{train}})^T \Sigma_{\text{train}}^{-1} (g_{\text{target}} - \mu_{\text{train}})}
\end{equation}
where $\mu_{\text{train}}$ and $\Sigma_{\text{train}}$ denote the empirical mean and regularized covariance of training hardware specifications.

\subsection{Component D: Group-Calibrated Conformal Risk Calibration}
To construct an empirical $(1-\alpha)$ safety bound under hardware shift, we employ a shift-adjusted conformal calibration framework. On an in-distribution calibration split, non-conformity scores are computed using total ensemble predictive uncertainty ($\hat{\sigma}_i = \hat{\sigma}_{\text{total}}(x_i, g_i)$) as normalized residuals $s_i = \frac{|y_i - \hat{T}_i|}{\hat{\sigma}_i}$. The base conformal quantile $\hat{q}_{1-\alpha}$ ($\alpha = 0.10$) is derived as:
\begin{equation}
\hat{q}_{1-\alpha} = \text{Quantile}_{\frac{\lceil (n+1)(1-\alpha) \rceil}{n}}\left(\{s_i\}_{i=1}^n\right)
\end{equation}

Under microarchitectural covariate shift, the effective calibration quantile is scaled proportionally to the clipped hardware discrepancy score ($\gamma_{\text{shift}} = 0.15, d_{\max} = 5.0$):
\begin{equation}
\hat{q}_{\text{shift}}(g) = \hat{q}_{1-\alpha} \cdot \left(1 + \gamma_{\text{shift}} \cdot \min\left(d_{\text{shift}}(g, \mathcal{D}_{\text{train}}), d_{\max}\right)\right)
\end{equation}

The empirical upper latency bound used for SLA verification in Eq.~\eqref{eq:opt_problem} is then formulated as:
\begin{equation}
\hat{T}_{\text{upper}}(x, g) = \hat{T}(x, g) + \hat{q}_{\text{shift}}(g) \cdot \hat{\sigma}_{\text{total}}(x, g)
\end{equation}

\subsection{Component E: Active Few-Shot Calibration and Boundary Acquisition Heuristic}
When introducing an unfamiliar GPU architecture, SHIFT-GPU supports optional active few-shot calibration within a minimal measurement budget $k \in \{1, 3, 5\}$.

Traditional active learning policies select samples maximizing predictive variance ($\arg\max \sigma^2(x)$). In systems infrastructure, variance on workloads that are far from the feasibility threshold ($T \ll T_{\max}$ or $T \gg T_{\max}$) may provide limited boundary discrimination. To explore threshold sensitivity, we formulate a boundary-focused acquisition heuristic, denoted \textbf{ROAA} (\textit{Regret-Oriented Active Acquisition}):
\begin{equation}
\label{eq:roaa}
\alpha_{\text{ROAA}}(x) = \hat{\sigma}_{\text{total}}(x, g) \cdot \exp\left( - \frac{(\hat{T}(x, g) - T_{\max})^2}{2 \hat{\sigma}_{\text{total}}^2(x, g)} \right) \cdot \left(\frac{\hat{P}(x, g)}{TDP(g)}\right)
\end{equation}

The first term reflects total uncertainty, the second term evaluates Gaussian proximity to the critical SLA decision boundary $T_{\max}$, and the third term incorporates normalized power impact. We emphasize that Eq.~\eqref{eq:roaa} is an empirical heuristic designed to sample boundary edge cases rather than a provably optimal policy; its objective is to explore boundary configurations rather than minimize global regression error across all workloads. 

Given $k$ probed measurements, PICAN adapts each ensemble member's representation via Feature-wise Linear Modulation (FiLM) \citep{perez2018film} applied to the latent space:
\begin{equation}
z_{\text{adapted}} = \gamma_k \odot z + \beta_k
\end{equation}
where modulation vectors $\gamma_k, \beta_k \in \mathbb{R}^d$ are optimized for 35 Adam steps with a ridge anchor toward identity ($10^{-2}(\|\gamma_k - \mathbf{1}\|^2 + \|\beta_k\|^2)$) without modifying the underlying physics-guided backbone weights. Algorithm \ref{alg:shift_gpu} details the complete operational procedure.

\begin{algorithm}[t]
\caption{SHIFT-GPU Admission Control, Conformal Risk Calibration, and Active Adaptation}
\label{alg:shift_gpu}
\begin{algorithmic}[1]
\REQUIRE Workload $x = (m, s)$, Candidates $\mathcal{G}$, Target SLA $T_{\max}$, Confidence $1-\alpha$, Budget $k$.
\ENSURE Optimal Hardware Accelerator $g^*$.
\FORALL{$g \in \mathcal{G}$}
    \STATE Compute deterministic Roofline floor $T_{\text{roof}}(x, g)$ via Eq.~\eqref{eq:roofline}.
    \STATE Evaluate microarchitectural Mahalanobis shift $d_{\text{shift}}(g, \mathcal{D}_{\text{train}})$ via Eq.~\eqref{eq:mahalanobis}.
    \IF{$g$ is unseen and profiling budget $k > 0$}
        \STATE Select $k$ calibration workloads $\{x_1, \dots, x_k\}$ maximizing $\alpha_{\text{ROAA}}(x)$ via Eq.~\eqref{eq:roaa}.
        \STATE Profile $\{x_j\}$ on $g$ to measure ground-truth $\{T_j, P_j\}$.
        \STATE Calibrate FiLM latent modulation parameters $[\gamma_k, \beta_k]$.
    \ENDIF
    \STATE Infer predicted mean $\hat{T}(x, g)$, power $\hat{P}(x, g)$, and log-variance $\hat{\sigma}^2(x, g)$ via PICAN.
    \STATE Compute empirical upper latency bound $\hat{T}_{\text{upper}}(x, g) = \hat{T}(x, g) + \hat{q}_{\text{shift}}(g) \hat{\sigma}(x, g)$.
\ENDFOR
\STATE Determine feasible set $\mathcal{G}_{\text{feasible}} = \{g \in \mathcal{G} \mid \hat{T}_{\text{upper}}(x, g) \le T_{\max}\}$.
\IF{$\mathcal{G}_{\text{feasible}} \neq \emptyset$}
    \STATE \textbf{return} $g^* = \arg\min_{g \in \mathcal{G}_{\text{feasible}}} \hat{P}(x, g)$
\ELSE
    \STATE \textbf{return} Abstain or Flag Cluster Capacity Deficit.
\ENDIF
\end{algorithmic}
\end{algorithm}

\section{Experimental Setup}
\label{sec:experiments}

\subsection{Dataset and Preprocessing}
We evaluate SHIFT-GPU using the public WattGPU benchmark repository \citep{fadel2026wattgpu}, derived from the large-scale Watt Counts telemetry corpus \citep{fadel2026wattcounts}. The dataset comprises 4,163 rigorously validated inference executions spanning 42 open-source autoregressive LLMs (ranging from 124M parameter GPT-2 to 27B parameter Gemma-3) and 8 server-grade NVIDIA accelerator architectures. All executions were measured on bare-metal production nodes running vLLM under controlled hardware telemetry sampling (NVIDIA NVML) across offline and Poisson-arrival server scenarios.

\subsection{Hardware Accelerator Population}
The evaluated accelerators span five distinct NVIDIA architectural generations, summarized in Table \ref{tab:hardware_specs}.

\begin{table}[h]
\centering
\caption{Microarchitectural specifications of evaluated server-grade GPU hardware population.}
\label{tab:hardware_specs}
\resizebox{\linewidth}{!}{%
\begin{tabular}{llrrrrrr}
\toprule
\textbf{GPU Model} & \textbf{Architecture} & \textbf{Release} & \textbf{Mem (GB)} & \textbf{BW (GB/s)} & \textbf{TFLOPS} & \textbf{TDP (W)} & \textbf{Process} \\
\midrule
Tesla V100 SXM2 & Volta & 2018 & 32.0 & 898.0 & 125.0 & 300 & 12 nm \\
Tesla T4 & Turing & 2018 & 16.0 & 320.0 & 65.0 & 70 & 12 nm \\
A100 SXM4 & Ampere & 2020 & 40.0 & 1560.0 & 312.0 & 400 & 7 nm \\
A30 PCIe & Ampere & 2021 & 24.0 & 933.1 & 165.0 & 165 & 7 nm \\
L40S & Ada Lovelace & 2022 & 48.0 & 864.0 & 362.0 & 350 & 5 nm \\
L4 & Ada Lovelace & 2023 & 24.0 & 300.1 & 121.0 & 72 & 5 nm \\
H100 NVL & Hopper & 2023 & 94.0 & 3940.0 & 835.0 & 400 & 5 nm \\
H200 NVL & Hopper & 2024 & 141.0 & 4890.0 & 835.0 & 700 & 5 nm \\
\bottomrule
\end{tabular}%
}
\end{table}

\subsection{Evaluation Protocols}
To rigorously assess generalization under out-of-distribution shifts, we execute two primary protocols:
\begin{enumerate}
    \item \textbf{Leave-One-GPU-Out (LOGO):} All models are trained on 7 GPU architectures and evaluated exclusively on the 8th held-out architecture. This process is repeated across all 8 folds to evaluate generalization to unfamiliar hardware.
    \item \textbf{Leave-One-LLM-Out (LOLO):} Evaluates generalization across unseen foundation model architectures by holding out distinct model families.
\end{enumerate}

\subsection{Strict Zero-Leakage Blinding Verification Protocol}
A critical vulnerability in out-of-distribution benchmark evaluation is inadvertent target leakage, where test measurements inform feature scaling, conformal non-conformity quantiles, or hardware candidate ranking. To ensure rigorous zero-leakage validation across all 8 LOGO folds:
\begin{enumerate}
    \item \textbf{Quarantined Partitioning:} For each held-out GPU $g$, the remaining 7 GPUs are partitioned into disjoint Training (80\%) and Calibration (20\%) sets. Feature scalers (\texttt{StandardScaler}) and predictive models are fitted strictly on Training data. Conformal non-conformity quantiles are calibrated strictly on Calibration data.
    \item \textbf{Target-Label Blinding Verification:} To empirically verify zero target leakage, all test target labels ($\text{ITL}$ and power) are wiped to $\text{NaN}$ in an isolated evaluation dataframe copy prior to preprocessing and inference. We verify that predictions, empirical latency bounds, and cluster recommendations generated under blinded inputs matched within numerical tolerance ($\max |\hat{y}_{\text{blinded}} - \hat{y}_{\text{eval}}| \le 10^{-7}$) against the evaluation pipeline.
    \item \textbf{Decision Locking:} Cluster recommendations and policy selections are locked prior to revealing test targets. SLA violation counts and regret are scored exclusively post-locking.
\end{enumerate}

\subsection{Evaluation Metrics}
Models are evaluated across three dimensions:
\begin{itemize}
    \item \textbf{Regression Quality:} Median Absolute Percentage Error (MdAPE \%), Pearson correlation coefficient ($r$), and Kendall rank correlation coefficients for GPU ordering ($\tau_{\text{GPU}}$) and LLM ordering ($\tau_{\text{LLM}}$).
    \item \textbf{Uncertainty Calibration:} Empirical Prediction Interval Coverage (\% at nominal 90\% level) and Mean Interval Width.
    \item \textbf{System Decision Quality:} Latency SLA Violation Rate (\%), Mean Selected Power (Watts), Operational Recommendation Regret (Eq.~\eqref{eq:regret}), and Cluster Abstention Frequency (\%).
\end{itemize}

\section{Experimental Results and Analysis}
\label{sec:results}

\subsection{E1 \& E2: Unseen Hardware (LOGO) and Unseen Model (LOLO) Generalization}
Table \ref{tab:logo_results} reports the performance of the 5-seed SHIFT-GPU (PICAN) ensemble against our matched XGBoost baseline rerun \citep{fadel2026wattgpu}, a Roofline-augmented XGBoost baseline predicting $\log(T / T_{\text{roof}})$, and pure physical bounds across all 8 held-out GPU architectures under LOGO cross-validation. Table \ref{tab:lolo_results} reports 42-fold Leave-One-LLM-Out (LOLO) generalization across unseen foundation models.

\begin{table}[h]
\centering
\caption{Zero-shot performance comparison on unseen hardware architectures under the Leave-One-GPU-Out (LOGO) protocol across 4,163 configurations. SHIFT-GPU reports 5-seed ensemble accuracy alongside single-seed mean $\pm$ std across $\{42, 123, 456, 789, 2026\}$. Lower MdAPE indicates superior accuracy; higher $r$ and $\tau$ indicate superior correlation and ranking.}
\label{tab:logo_results}
\resizebox{\linewidth}{!}{%
\begin{tabular}{llrrrr}
\toprule
\textbf{Target Metric} & \textbf{Model Architecture} & \textbf{MdAPE (\%)} & \textbf{Pearson $r$} & \textbf{GPU $\tau$} & \textbf{LLM $\tau$} \\
\midrule
\multirow{4}{*}{\textbf{Inter-Token Latency (ITL)}} 
& Non-Learning Roofline \citep{williams2009roofline} & 38.10 & 0.992 & \textbf{0.822} & 0.908 \\
& Matched XGBoost Rerun \citep{fadel2026wattgpu} & 11.98 & 0.983 & 0.676 & 0.922 \\
& XGBoost + Roofline Baseline & 10.68 & 0.992 & 0.687 & 0.933 \\
& \textbf{SHIFT-GPU (5-Seed PICAN Ensemble)} & \textbf{6.83} (7.22$\pm$0.75) & \textbf{0.995} & 0.757 & \textbf{0.943} \\
\midrule
\multirow{3}{*}{\textbf{Mean Power Draw (W)}} 
& Non-Learning TDP Duty Base & 27.35 & 0.834 & 0.616 & 0.768 \\
& \textbf{Matched XGBoost Rerun} \citep{fadel2026wattgpu} & \textbf{12.52} & 0.973 & 0.681 & \textbf{0.812} \\
& SHIFT-GPU (5-Seed PICAN Ensemble) & 13.44 (14.00$\pm$1.56) & \textbf{0.978} & \textbf{0.707} & \textbf{0.812} \\
\bottomrule
\end{tabular}%
}
\end{table}

\begin{table}[h]
\centering
\caption{Cross-model Leave-One-LLM-Out (LOLO) ITL generalization across 42 held-out LLMs.}
\label{tab:lolo_results}
\resizebox{0.82\linewidth}{!}{%
\begin{tabular}{lrrrr}
\toprule
\textbf{Model Architecture} & \textbf{MdAPE (\%)} & \textbf{Pearson $r$} & \textbf{GPU $\tau$} & \textbf{LLM $\tau$} \\
\midrule
Matched XGBoost Rerun \citep{fadel2026wattgpu} & 7.72 & 0.992 & 0.874 & 0.903 \\
XGBoost + Roofline Baseline & 6.72 & \textbf{0.996} & 0.881 & 0.906 \\
\textbf{SHIFT-GPU (PICAN)} & \textbf{6.21} & 0.995 & \textbf{0.885} & \textbf{0.917} \\
\bottomrule
\end{tabular}%
}
\end{table}

As shown in Table \ref{tab:logo_results}, PICAN reduces unseen-GPU ITL prediction error from 11.98\% (matched XGBoost rerun) and 10.68\% (Roofline-augmented XGBoost) down to \textbf{6.83\%} (a 43.0\% relative error reduction against WattGPU XGBoost and 36.1\% against XGBoost + Roofline), with consistent single-seed stability ($7.22\% \pm 0.75\%$ across 5 seeds) and top LLM ranking fidelity ($\tau = 0.943$). On unseen model architectures (42-fold LOLO, Table \ref{tab:lolo_results}), PICAN similarly reduces ITL error to \textbf{6.21\%} vs.\ 7.72\% for XGBoost ($-$19.5\%) and 6.72\% for XGBoost + Roofline ($-$7.6\%). For power draw under LOGO, PICAN achieves 13.44\% MdAPE (vs.\ 27.35\% duty-cycle baseline and 12.52\% for normalized-power XGBoost), attaining higher linear correlation ($r = 0.978$) and cross-GPU rank ordering ($\tau = 0.707$ vs.\ 0.681) while XGBoost edges out on median percentage error because static TDP bounds lack an exact dynamic utilization floor analogous to Roofline for ITL.

\begin{table}[h]
\centering
\caption{Per-GPU zero-shot Inter-Token Latency prediction error across all 8 held-out GPU architectures under LOGO cross-validation, comparing Matched XGBoost, XGBoost + Roofline, and 5-seed SHIFT-GPU (PICAN).}
\label{tab:per_gpu_results}
\resizebox{\linewidth}{!}{%
\begin{tabular}{llrrrrrrr}
\toprule
\textbf{Held-Out GPU} & \textbf{Arch.} & \textbf{Shift $d_{\text{shift}}$} & \textbf{XGBoost} & \textbf{XGB+Roof} & \textbf{SHIFT-GPU} & \textbf{Red.\ vs.\ XGB} & \textbf{2-Sided Cov.} & \textbf{Upper Cov.} \\
\midrule
A100 SXM4 40 GB & Ampere & 1.42 & \textbf{8.95\%} & 12.27\% & 11.87\% & $-$32.6\% & 62.8\% & 100.0\% \\
A30 PCIe & Ampere & 3.41 & 17.15\% & 5.33\% & \textbf{5.03\%} & +70.7\% & 82.9\% & 82.9\% \\
H100 NVL 94 GB & Hopper & 6.25 & 8.36\% & 7.18\% & \textbf{4.49\%} & +46.2\% & 65.4\% & 65.6\% \\
H200 NVL & Hopper & 17.03 & 16.73\% & 7.08\% & \textbf{6.55\%} & +60.9\% & 75.0\% & 100.0\% \\
L4 & Ada & 4.95 & 17.39\% & 15.98\% & \textbf{3.31\%} & +81.0\% & 99.1\% & 100.0\% \\
L40S & Ada & 10.61 & 10.21\% & 7.20\% & \textbf{4.50\%} & +55.9\% & 98.5\% & 98.7\% \\
Tesla T4 & Turing & 2.55 & 13.65\% & \textbf{11.18\%} & 13.42\% & +1.7\% & 73.2\% & 95.1\% \\
Tesla V100 SXM2 & Volta & 4.30 & 11.74\% & 16.89\% & \textbf{10.21\%} & +13.0\% & 68.8\% & 68.8\% \\
\midrule
\textbf{Overall Population} & \textbf{8 GPUs} & \textbf{4.95$^\dagger$} & \textbf{11.98\%} & \textbf{10.68\%} & \textbf{6.83\%} & \textbf{+43.0\%} & \textbf{77.0\%} & \textbf{86.8\%} \\
\bottomrule
\multicolumn{9}{l}{\footnotesize $^\dagger$Median shift distance across all 4,163 evaluation samples (GPU-level mean: 6.32). Rightmost columns report conformal coverage.}
\end{tabular}%
}
\end{table}

The disaggregated breakdown in Table \ref{tab:per_gpu_results} shows that SHIFT-GPU improves upon Roofline-augmented XGBoost on 7 of 8 GPUs and achieves low zero-shot error on modern data center GPUs: \textbf{3.31\%} on L4 (vs.\ 17.39\% XGBoost, 15.98\% XGB+Roofline), \textbf{4.49\%} on H100 NVL, \textbf{4.50\%} on L40S, \textbf{5.03\%} on A30 PCIe, and \textbf{6.55\%} on H200 NVL ($d_{\text{shift}} = 17.03$). Figure \ref{fig:temporal_dynamics}(a) illustrates parity across all 4,163 configurations, while Figure \ref{fig:temporal_dynamics}(b) shows the empirical cumulative error distributions.

\begin{figure*}[t]
\centering
\includegraphics[width=\textwidth]{figures/fig2_temporal_dynamics_panel.png}
\caption{Zero-shot prediction calibration and error distributions across 4,163 empirical configurations. (a) Parity plot comparing measured vs.\ predicted Inter-Token Latency ($T_{\mathrm{meas}}$ vs.\ $\hat{T}$) on a log--log scale for SHIFT-GPU ($r = 0.995$, MdAPE = 6.83\%) and Matched XGBoost ($r = 0.983$, MdAPE = 11.98\%). (b) Empirical Cumulative Distribution Function (CDF) of relative prediction error comparing SHIFT-GPU (6.83\%), XGBoost + Roofline (10.68\%), Matched XGBoost (11.98\%), and Non-Learning Roofline (38.10\%).}
\label{fig:temporal_dynamics}
\end{figure*}

\subsection{E3 \& E4: Microarchitectural Covariate Shift and Conformal Calibration}
We evaluate prediction interval validity under hardware covariate shift. Under strict LOGO cross-validation, uncalibrated Gaussian intervals formed from our trained heteroscedastic + ensemble standard deviation ($\hat{T} \pm 1.645\,\hat{\sigma}_{\text{total}}$) attain empirical coverage of \textbf{69.5\%} two-sided (\textbf{82.8\%} upper-bound) against the nominal 90\% target, confirming that hardware domain shift causes under-coverage.

In contrast, our shift-adjusted conformal calibrator scales calibration quantiles proportionally to clipped Mahalanobis shift distance ($\gamma_{\text{shift}}=0.15, d_{\max}=5.0$), lifting overall two-sided empirical coverage to \textbf{77.0\%} and upper-bound SLA coverage to \textbf{86.8\%} (reaching 95.1\%--100.0\% on A100, H200, L4, L40S, and T4). Empirical coverage remains below nominal 90.0\% on H100 (65.6\%) and V100 (68.8\%) because distinct GPU populations violate strict exchangeability. Figure \ref{fig:accuracy_efficiency}(a) and (b) plot per-GPU LOGO errors across all three baselines ordered by $d_{\text{shift}}$.

\begin{figure*}[t]
\centering
\includegraphics[width=\textwidth]{figures/fig3_accuracy_efficiency_panel.png}
\caption{Hardware generalization across unseen accelerator targets under microarchitectural covariate shift. (a) Leave-one-group-out (LOGO) ITL prediction MdAPE across 8 GPU architectures ordered by Mahalanobis shift distance $d_{\mathrm{shift}}$, comparing Matched XGBoost, XGBoost + Roofline, and SHIFT-GPU. (b) Error scaling trends as a function of $d_{\mathrm{shift}}$.}
\label{fig:accuracy_efficiency}
\end{figure*}

\subsection{E5: Active Few-Shot Hardware Profiling Benchmark}
Table \ref{tab:few_shot_results} and Figure \ref{fig:performance_panel} evaluate few-shot FiLM adaptation across profiling budgets $k \in \{0, 1, 3, 5\}$ and four acquisition strategies strictly on untouched test workloads ($D_{\text{test}} \setminus D_{\text{probe}}$), driving uncertainty and ROAA selection with the trained total ensemble standard deviation $\hat{\sigma}_{\text{total}}$.

\begin{table}[h]
\centering
\caption{Few-shot FiLM adaptation benchmark on unseen accelerators across measurement budgets ($k$) and acquisition policies. Evaluated strictly on untouched test samples ($D_{\text{test}} \setminus D_{\text{probe}}$).}
\label{tab:few_shot_results}
\resizebox{0.88\linewidth}{!}{%
\begin{tabular}{llrr}
\toprule
\textbf{Profiling Budget ($k$)} & \textbf{Acquisition Policy} & \textbf{MdAPE (\%)} & \textbf{Pearson $r$} \\
\midrule
$k = 0$ (Zero-Shot Ensemble) & Untrained Target Baseline & \textbf{6.83} & 0.995 \\
\midrule
\multirow{4}{*}{$k = 1$ (One-Shot)} 
& Uniform Random & 7.62 & 0.997 \\
& K-Means Feature Diversity & 10.13 & 0.997 \\
& Maximal Predictive Uncertainty & 7.86 & 0.997 \\
& Boundary Heuristic (ROAA) & 7.93 & 0.996 \\
\midrule
\multirow{4}{*}{$k = 3$ (Three-Shot)} 
& Uniform Random & 7.37 & 0.997 \\
& K-Means Feature Diversity & 7.78 & 0.997 \\
& Maximal Predictive Uncertainty & 7.95 & 0.997 \\
& Boundary Heuristic (ROAA) & 8.00 & 0.996 \\
\midrule
\multirow{4}{*}{$k = 5$ (Five-Shot)} 
& Uniform Random & 7.38 & \textbf{0.997} \\
& K-Means Feature Diversity & 7.56 & \textbf{0.997} \\
& Maximal Predictive Uncertainty & 7.98 & \textbf{0.997} \\
& Boundary Heuristic (ROAA) & 8.10 & 0.996 \\
\bottomrule
\end{tabular}%
}
\end{table}

Two clear empirical insights emerge from Table \ref{tab:few_shot_results}. First, because the 5-seed zero-shot ensemble already generalizes strongly across unseen hardware (\textbf{6.83\%} MdAPE), adapting the 128 FiLM parameters from only $k \in \{1, 3, 5\}$ shots slightly increases global MdAPE ($7.37\%$--$8.10\%$) while tightening tail linearity (Pearson $r$ improves from 0.995 to 0.997). Second, consistent across all tested budgets ($k \in \{1, 3, 5\}$), Uniform Random sampling outperforms K-Means diversity, maximal uncertainty, and the boundary-targeted ROAA heuristic ($7.37\%$--$7.62\%$ vs.\ $7.93\%$--$8.10\%$). Because ROAA concentrates shots tightly around a single latency threshold $T_{\max} = 35$~ms, updating global modulation vectors on boundary-skewed probes distorts predictions across the broader workload distribution. Consequently, we advise practitioners to deploy the zero-shot ensemble directly or use uniform random probing when larger calibration sets become available. Figure \ref{fig:performance_panel}(b) confirms smooth scaling across all 42 dense LLMs spanning \textbf{0.12B to 27B parameters}.

\begin{figure*}[t]
\centering
\includegraphics[width=\textwidth]{figures/fig4_computational_performance_panel.png}
\caption{Few-shot FiLM adaptation and architectural parameter scalability. (a) Inter-token latency MdAPE across budgets $k \in \{0, 1, 3, 5\}$ comparing Uniform Random, K-Means Diversity, Max Uncertainty, and Boundary Heuristic (ROAA). (b) Architectural scalability comparing measured median ITL vs.\ SHIFT-GPU predictions across 42 evaluated LLMs spanning 0.12B to 27B parameters.}
\label{fig:performance_panel}
\end{figure*}

\subsection{E6: Production Hardware Recommendation and SLA Regret}
We simulate a production cloud scheduler serving diverse LLM requests under an SLA latency threshold $T_{\max} = 35$ ms. Table \ref{tab:recommendations} and Figures \ref{fig:per_architecture_error}--\ref{fig:regret_panel} compare seven operational selection policies evaluated on the same 125 cluster workloads using Eq. (\ref{eq:regret}) regret.

\begin{table}[h]
\centering
\caption{Production hardware recommendation benchmark under 35 ms latency SLA. Evaluated across 125 unique LLM workloads over 8 candidate GPU architectures using Eq. (\ref{eq:regret}) regret ($\lambda = 5$).}
\label{tab:recommendations}
\resizebox{\linewidth}{!}{%
\begin{tabular}{lrrrrr}
\toprule
\textbf{Selection Policy} & \textbf{Accepted Workloads} & \textbf{Abstention (\%)} & \textbf{SLA Violations} & \textbf{Mean Power (W)} & \textbf{Mean Regret} \\
\midrule
Minimum TDP (Naive) & 125 / 125 & 0.0 & 44 / 125 (35.2\%) & \textbf{102.5} & 142.9 \\
Minimum TDP + Roofline & 125 / 125 & 0.0 & 12 / 125 (9.6\%) & 106.3 & 65.4 \\
Point Estimate (Matched XGBoost) & 125 / 125 & 0.0 & 0 / 125 (0.0\%) & 105.7 & \textbf{2.3} \\
Point Estimate (XGBoost + Roofline) & 125 / 125 & 0.0 & 1 / 125 (0.8\%) & 105.7 & 4.7 \\
Point Estimate (SHIFT-GPU PICAN) & 125 / 125 & 0.0 & 0 / 125 (0.0\%) & 115.2 & 11.8 \\
\textbf{SHIFT-GPU (Risk-Calibrated)} & \textbf{125 / 125} & \textbf{0.0} & \textbf{0 / 125 (0.0\%)} & 115.0 & 11.6 \\
True Empirical Oracle & 125 / 125 & 0.0 & 0 / 125 (0.0\%) & 103.4 & 0.0 \\
\bottomrule
\end{tabular}%
}
\end{table}

\begin{table}[h]
\centering
\caption{Multi-SLA sensitivity benchmark evaluating cluster selection policies across latency thresholds $T_{\max} \in \{20, 25, 30, 35, 40\}$~ms on 125 unique LLM workloads.}
\label{tab:sla_sensitivity}
\resizebox{\linewidth}{!}{%
\begin{tabular}{llrrrrr}
\toprule
\textbf{Latency SLA ($T_{\max}$)} & \textbf{Policy} & \textbf{Accepted Workloads} & \textbf{Abstain (\%)} & \textbf{SLA Violations} & \textbf{Mean Power (W)} & \textbf{Mean Regret} \\
\midrule
\multirow{7}{*}{$T_{\max} = 20$~ms} 
& Minimum TDP (Naive) & 125 / 125 & 0.0 & 64 / 125 (51.2\%) & 102.5 & 185.5 \\
& Minimum TDP + Roofline & 125 / 125 & 0.0 & 31 / 125 (24.8\%) & 104.8 & 98.1 \\
& Matched XGBoost Point & 125 / 125 & 0.0 & 2 / 125 (1.6\%) & 118.4 & 12.2 \\
& XGBoost + Roofline Point & 125 / 125 & 0.0 & 0 / 125 (0.0\%) & 113.3 & 3.6 \\
& SHIFT-GPU PICAN Point & 125 / 125 & 0.0 & 3 / 125 (2.4\%) & 118.1 & 15.0 \\
& \textbf{SHIFT-GPU (Risk-Calibrated)} & \textbf{125 / 125} & \textbf{0.0} & \textbf{0 / 125 (0.0\%)} & \textbf{121.1} & \textbf{11.4} \\
& True Empirical Oracle & 125 / 125 & 0.0 & 0 / 125 (0.0\%) & 109.6 & 0.0 \\
\midrule
\multirow{7}{*}{$T_{\max} = 25$~ms} 
& Minimum TDP (Naive) & 125 / 125 & 0.0 & 57 / 125 (45.6\%) & 102.5 & 176.0 \\
& Minimum TDP + Roofline & 125 / 125 & 0.0 & 11 / 125 (8.8\%) & 103.7 & 20.9 \\
& Matched XGBoost Point & 125 / 125 & 0.0 & 3 / 125 (2.4\%) & 109.5 & 10.8 \\
& XGBoost + Roofline Point & 125 / 125 & 0.0 & 1 / 125 (0.8\%) & 107.0 & 4.5 \\
& SHIFT-GPU PICAN Point & 125 / 125 & 0.0 & 3 / 125 (2.4\%) & 116.3 & 16.5 \\
& \textbf{SHIFT-GPU (Risk-Calibrated)} & \textbf{125 / 125} & \textbf{0.0} & \textbf{0 / 125 (0.0\%)} & \textbf{117.1} & \textbf{12.3} \\
& True Empirical Oracle & 125 / 125 & 0.0 & 0 / 125 (0.0\%) & 104.9 & 0.0 \\
\midrule
\multirow{7}{*}{$T_{\max} = 30$~ms} 
& Minimum TDP (Naive) & 125 / 125 & 0.0 & 49 / 125 (39.2\%) & 102.5 & 161.1 \\
& Minimum TDP + Roofline & 125 / 125 & 0.0 & 7 / 125 (5.6\%) & 104.6 & 29.1 \\
& Matched XGBoost Point & 125 / 125 & 0.0 & 1 / 125 (0.8\%) & 106.3 & 4.8 \\
& XGBoost + Roofline Point & 125 / 125 & 0.0 & 1 / 125 (0.8\%) & 105.9 & 4.5 \\
& SHIFT-GPU PICAN Point & 125 / 125 & 0.0 & 2 / 125 (1.6\%) & 115.6 & 15.4 \\
& \textbf{SHIFT-GPU (Risk-Calibrated)} & \textbf{125 / 125} & \textbf{0.0} & \textbf{0 / 125 (0.0\%)} & \textbf{116.1} & \textbf{12.3} \\
& True Empirical Oracle & 125 / 125 & 0.0 & 0 / 125 (0.0\%) & 103.8 & 0.0 \\
\midrule
\multirow{7}{*}{$T_{\max} = 35$~ms} 
& Minimum TDP (Naive) & 125 / 125 & 0.0 & 44 / 125 (35.2\%) & 102.5 & 142.9 \\
& Minimum TDP + Roofline & 125 / 125 & 0.0 & 12 / 125 (9.6\%) & 106.3 & 65.4 \\
& Matched XGBoost Point & 125 / 125 & 0.0 & 0 / 125 (0.0\%) & 105.7 & 2.3 \\
& XGBoost + Roofline Point & 125 / 125 & 0.0 & 1 / 125 (0.8\%) & 105.7 & 4.7 \\
& SHIFT-GPU PICAN Point & 125 / 125 & 0.0 & 0 / 125 (0.0\%) & 115.2 & 11.8 \\
& \textbf{SHIFT-GPU (Risk-Calibrated)} & \textbf{125 / 125} & \textbf{0.0} & \textbf{0 / 125 (0.0\%)} & \textbf{115.0} & \textbf{11.6} \\
& True Empirical Oracle & 125 / 125 & 0.0 & 0 / 125 (0.0\%) & 103.4 & 0.0 \\
\midrule
\multirow{7}{*}{$T_{\max} = 40$~ms} 
& Minimum TDP (Naive) & 125 / 125 & 0.0 & 39 / 125 (31.2\%) & 102.5 & 119.2 \\
& Minimum TDP + Roofline & 125 / 125 & 0.0 & 11 / 125 (8.8\%) & 105.7 & 49.6 \\
& Matched XGBoost Point & 125 / 125 & 0.0 & 0 / 125 (0.0\%) & 105.5 & 2.2 \\
& XGBoost + Roofline Point & 125 / 125 & 0.0 & 0 / 125 (0.0\%) & 105.5 & 2.1 \\
& SHIFT-GPU PICAN Point & 125 / 125 & 0.0 & 2 / 125 (1.6\%) & 114.1 & 16.0 \\
& \textbf{SHIFT-GPU (Risk-Calibrated)} & \textbf{125 / 125} & \textbf{0.0} & \textbf{0 / 125 (0.0\%)} & \textbf{115.2} & \textbf{11.8} \\
& True Empirical Oracle & 125 / 125 & 0.0 & 0 / 125 (0.0\%) & 103.4 & 0.0 \\
\bottomrule
\end{tabular}%
}
\end{table}

As summarized in Table \ref{tab:recommendations}, under $T_{\max} = 35$~ms, non-learning heuristics fail severely: naive Minimum TDP incurs 44 / 125 SLA violations (35.2\%, mean regret 142.9), while augmenting Minimum TDP with the deterministic Roofline filter still leaves 12 / 125 violations (9.6\%, mean regret 65.4). Both SHIFT-GPU point estimation (115.2~W) and risk-calibrated SHIFT-GPU (115.0~W) record 0 / 125 violations (0.0\%), whereas XGBoost + Roofline point selection incurs 1 / 125 violation (0.8\%). Matched XGBoost point prediction achieves 0 violations specifically at 35~ms with 105.7~W due to slightly tighter zero-shot power estimation (12.52\% vs.\ 13.44\% Power MdAPE), which ranks lower-TDP feasible candidates closer to the True Oracle (103.4~W).

To examine robustness under tightening constraints, Table \ref{tab:sla_sensitivity} sweeps $T_{\max} \in \{20, 25, 30, 35, 40\}$~ms. Across the stricter thresholds, every uncalibrated point policy breaches SLA bounds: Matched XGBoost incurs 2 violations (1.6\%) at 20~ms, 3 (2.4\%) at 25~ms, and 1 (0.8\%) at 30~ms; XGBoost + Roofline incurs 1 violation (0.8\%) at 25, 30, and 35~ms; and SHIFT-GPU PICAN Point incurs 2--3 violations (1.6\%--2.4\%) at 20, 25, 30, and 40~ms. Crucially, \textbf{risk-calibrated SHIFT-GPU is the sole policy to maintain 0 / 125 (0.0\%) SLA violations with 0 abstentions across all five SLA thresholds}, adding a modest risk-hedging power margin (+3.0~W at 20~ms, +0.8~W at 25~ms, +0.5~W at 30~ms relative to PICAN Point) while lower-bounding cumulative operational regret below Matched XGBoost under the tightest 20~ms constraint (mean regret 11.4 vs.\ 12.2; Figure \ref{fig:regret_panel}(a)).

\begin{figure*}[t]
\centering
\includegraphics[width=\textwidth]{figures/fig5_per_architecture_error_panel.png}
\caption{Empirical relative error distributions across NVIDIA GPU microarchitectural families on a log scale comparing Matched XGBoost, XGBoost + Roofline, and SHIFT-GPU (5-seed PICAN): (a) Ampere (A100 and A30), (b) Ada Lovelace (L40S and L4), and (c) Hopper (H100 and H200). Diamonds denote distribution means.}
\label{fig:per_architecture_error}
\end{figure*}

\begin{figure*}[t]
\centering
\includegraphics[width=\textwidth]{figures/fig6_cumulative_regret_panel.png}
\caption{Cumulative operational deployment regret under Eq. (\ref{eq:regret}) over 125 evaluation workloads for (a) $T_{\max} = 20$~ms, (b) $T_{\max} = 30$~ms, and (c) $T_{\max} = 35$~ms. Uncalibrated Minimum TDP (Roofline) accumulates severe SLA violation penalties, while uncalibrated point predictors exhibit intermittent regret jumps from SLA breaches; Risk-Calibrated SHIFT-GPU maintains a smooth, violation-free trajectory across all thresholds.}
\label{fig:regret_panel}
\end{figure*}

\subsection{E7: Systematic Ablation Study}
Table \ref{tab:ablations} and Figure \ref{fig:summary_heatmap} detail systematic ablations isolating the architectural and physics-informed contributions of SHIFT-GPU, retrained independently across the identical 5-seed LOGO protocol.

\begin{table}[h]
\centering
\caption{Systematic ablation study isolating architectural and algorithmic components of SHIFT-GPU. All neural variants trained with the identical 5-seed Gaussian NLL protocol across the 8 LOGO folds.}
\label{tab:ablations}
\resizebox{\linewidth}{!}{%
\begin{tabular}{lrrr}
\toprule
\textbf{Model Architecture / Variant} & \textbf{ITL MdAPE (\%)} & \textbf{Upper Coverage (\%)} & \textbf{SLA Violation Rate (35 ms)} \\
\midrule
\textbf{Full Proposed SHIFT-GPU (PICAN + Roofline + Conformal)} & \textbf{6.83} & \textbf{86.8} & \textbf{0 / 125 (0.0\%)} \\
Ablation A: Without Physics Roofline (Pure Attention) & 31.70 & 90.9 & 5 / 125 (4.0\%) \\
Ablation B: Without Cross-Attention (Flat MLP + Roofline) & 6.42 & 87.4 & 1 / 125 (0.8\%) \\
Baseline C: XGBoost + Roofline & 10.68 & 83.8 & 1 / 125 (0.8\%) \\
Baseline D: Matched XGBoost \citep{fadel2026wattgpu} & 11.98 & 84.5 & 0 / 125 (0.0\%) \\
\bottomrule
\end{tabular}%
}
\end{table}

Ablating the deterministic microarchitectural Roofline base (Ablation A) causes zero-shot ITL prediction error to surge from 6.83\% to 31.70\% (+364\% relative error increase, $4.64\times$) and elevates SLA violations to 4.0\% (5 / 125), confirming that physical inductive anchoring is essential to prevent unconstrained neural models from extrapolating unphysical latencies on unfamiliar hardware. Replacing the bipartite cross-attention fusion with a flat concatenated MLP over the same Roofline prior and NLL training (Ablation B) yields 6.42\% MdAPE (slightly lower median error than PICAN's 6.83\%, well within the single-seed spread $\pm 0.75\%$), while incurring 1 / 125 (0.8\%) SLA violation compared to 0 / 125 (0.0\%) for SHIFT-GPU. As proved in Section \ref{sec:method}, because pairwise training uses singleton hardware tokens ($L_k = 1$), softmax attention weights equal $1.0$; consequently, both continuous neural Roofline-residual models generalize markedly better than axis-aligned trees, cutting LOGO error by 36.1\%--39.9\% relative to XGBoost + Roofline (10.68\% MdAPE, 0.8\% violations) and by 43.0\%--46.4\% relative to Matched XGBoost (11.98\% MdAPE). This directly isolates continuous neural Roofline-residual modeling as the key enabler of cross-GPU extrapolation.

\begin{figure*}[t]
\centering
\includegraphics[width=\textwidth]{figures/fig7_summary_heatmap.png}
\caption{Systematic component ablation and production recommendation trade-offs. (a) Zero-shot LOGO ITL MdAPE (\%) across five variants, showing that removing the physics Roofline base (Ablation A) surges error to 31.70\%, whereas both neural Roofline-residual models (6.83\% and 6.42\%) substantially outperform XGBoost + Roofline (10.68\%) and Matched XGBoost (11.98\%). (b) Production hardware selection frontier comparing mean cluster power draw (Watts) against SLA violation rate (\%) at $T_{\max} = 35$~ms.}
\label{fig:summary_heatmap}
\end{figure*}

\section{Discussion and Systems Implications}
\label{sec:discussion}

\subsection{Operational Energy vs. Instantaneous Power}
A common misconception in AI infrastructure is treating instantaneous GPU power draw (Watts) as synonymous with energy consumption (Joules). In autoregressive inference, total energy per request is the time integral of power: $E = \int_0^T P(t) dt \approx P_{\text{mean}} \cdot T_{\text{total}}$. A low-TDP accelerator (e.g., Tesla T4 at 70 W) operating at high latency may consume significantly more total energy than an advanced accelerator (e.g., L40S at 350 W) that finishes the inference burst in a fraction of the time. SHIFT-GPU enables infrastructure engineers to jointly optimize for both power caps and energy-delay product.

\subsection{Integration into Cloud Schedulers (Kubernetes / Slurm)}
SHIFT-GPU can be seamlessly integrated into existing cloud resource orchestrators as a lightweight admission controller, complementing node-level power coordination policies such as PowerCoord \citep{suscom_powercoord2020} and cloud workflow managers \citep{suscom_kumar2022, suscom_malekloo2018}. Furthermore, by explicitly bounding tail latency risk through conformal calibration, SHIFT-GPU guards against the queue instability and latency explosions analytically established by Yao et al. \citep{suscom_tradeoffs2025}. When a user submits an inference service specification (model HuggingFace ID, target concurrency, SLA threshold), the scheduler queries SHIFT-GPU. If the target GPU family is well covered by historical telemetry, SHIFT-GPU returns zero-shot recommendations in under 5 milliseconds using the 5-seed ensemble. When unfamiliar accelerators are introduced ($d_{\text{shift}} > 3.0$), the scheduler automatically widens conformal admission margins via Mahalanobis scaling and can optionally run a uniform random probing micro-job to log empirical target behavior once larger calibration batches are available.

\section{Threats to Validity and Limitations}
\label{sec:threats}

\textbf{Calibration under hardware shift:}
Split conformal prediction provides finite-sample marginal coverage of at least $1-\alpha$ under exchangeability of calibration and test examples. In our leave-one-GPU-out evaluation, calibration and test samples come from distinct GPU populations, so strict exchangeability cannot be assumed. Relative to uncalibrated Gaussian intervals (69.5\% two-sided, 82.8\% upper), shift-adjusted conformal calibration raises measured coverage to 86.8\% for the upper latency bound and 77.0\% for the two-sided interval, both still slightly below the nominal 90\% target due to extreme domain shifts on architectures such as H100 NVL (65.6\%) and V100 (68.8\%). The Mahalanobis adjustment is an empirical heuristic, and establishing finite-sample coverage guarantees under non-exchangeable hardware shift remains open.

\textbf{Workload-hardware attention dynamics:}
In our pairwise training and inference formulation, each input sample comprises a single workload specification $x$ and a single candidate accelerator $g$. Because the key sequence length is unity ($L_k = 1$), the softmax attention distribution is mathematically degenerate ($\mathrm{softmax}(s) = 1.0$), rendering query-key dot products non-operational for single-pair inference. In Ablation B, both decoupled PICAN (6.83\% MdAPE, 0.0\% SLA violations at 35~ms) and the flat concatenated MLP over the same Roofline prior (6.42\% MdAPE, 0.8\% SLA violations) fall within single-seed variation ($\pm 0.75\%$), whereas both dramatically outperform Roofline-augmented XGBoost (10.68\%) and raw XGBoost (11.98\%). The primary source of cross-GPU accuracy is therefore smooth neural Roofline-residual modeling rather than single-token cross-attention. Extending PICAN to joint list-wise multi-GPU encoding---where a workload query attends simultaneously over $N > 1$ candidate accelerators---remains a promising direction to unlock data-dependent attention routing.

\textbf{Power prediction and conservative hedging:}
Zero-shot power prediction exhibits higher relative variance on unfamiliar GPUs than latency because TDP scaling interacts non-linearly with per-SM boost clocks and memory utilization. Across LOGO evaluation, 5-seed PICAN records 13.44\% Power MdAPE compared to 12.52\% for Matched XGBoost (Table \ref{tab:logo_results}), because tree splits closely track static TDP regimes. Consequently, under relaxed 35~ms constraints, Matched XGBoost selects slightly lower-power configurations (105.7~W vs.\ 115.0~W), while risk-calibrated SHIFT-GPU trades this conservative +9.3~W margin for universal zero-violation reliability (0 / 125 across all SLA bounds from 20~ms to 40~ms). Jointly calibrating power bounds or adopting physics-informed thermal models for power residuals warrants future study.

\textbf{Few-shot active acquisition:}
While we formulated a boundary-targeted acquisition heuristic (ROAA) to probe near-threshold configurations, empirical results showed that uniform random sampling achieved lower global regression error across all tested budgets ($k \in \{1, 3, 5\}$; 7.37\%--7.62\% vs.\ 7.93\%--8.10\% MdAPE), and the zero-shot 5-seed ensemble already achieved lower global median error (6.83\%). Updating FiLM affine adapters on a handful of non-uniform or boundary-skewed shots perturbs global feature calibration across the broader workload space. Identifying acquisition criteria and regularization schedules that improve global adaptation at ultra-low shot counts remains an open problem.

\textbf{Hardware scope:}
The benchmark covers eight NVIDIA GPU models. Holding out one GPU model can leave other members of its architecture family in training. Generalization to unseen architecture families and other accelerator vendors remains to be evaluated.

\textbf{Software stack:}
The benchmark uses vLLM with FP16/BF16 precision. External validity across serving runtimes and lower-precision formats remains to be evaluated.

\section{Conclusion}
\label{sec:conclusion}

We presented SHIFT-GPU, a physics-guided, risk-calibrated framework for predicting LLM inference latency and selecting energy-efficient accelerators under hardware covariate shift. Under strict NaN-blinded leave-one-GPU-out evaluation across 4,163 benchmark configurations spanning 42 LLMs (0.12B--27B parameters) and 8 NVIDIA GPUs, SHIFT-GPU reduced inter-token latency MdAPE from 11.98\% (matched XGBoost) and 10.68\% (Roofline-augmented XGBoost) down to \textbf{6.83\%} (43.0\% and 36.1\% relative reductions), and achieved \textbf{6.21\% MdAPE} under 42-model leave-one-LLM-out evaluation. Systematic ablations isolate smooth neural Roofline-residual modeling as the central driver of unseen-GPU generalization, showing a 4.64$\times$ error surge to 31.70\% when removing the Roofline prior and substantial improvements over tree-based Roofline residuals.

In production cluster selection across five latency SLAs ($T_{\max} \in \{20, 25, 30, 35, 40\}$~ms), risk-calibrated SHIFT-GPU was the only evaluated policy to achieve \textbf{0 violations across all 125 workloads (0.0\%, 0 abstentions)} at every threshold (mean power 115.0~W at 35~ms and 121.1~W at 20~ms), whereas naive Minimum TDP incurred 31.2\%--51.2\% SLA violations, Roofline Minimum TDP incurred 5.6\%--24.8\%, and uncalibrated point predictors incurred up to 2.4\% violations under strict constraints. Combined with our transparent empirical characterization of $L_k=1$ attention equivalence and few-shot FiLM probing, SHIFT-GPU provides a rigorously validated, reproducible foundation for sustainable LLM infrastructure scheduling.

\section*{CRediT Authorship Contribution Statement}
\textbf{Jason Pandian:} Conceptualization, Methodology, Software, Validation, Formal analysis, Investigation, Data Curation, Writing - Original Draft, Writing - Review \& Editing, Visualization, Supervision, Project administration, Resources.

\section*{Declaration of Competing Interest}
The author declares that he has no known competing financial interests or personal relationships that could have appeared to influence the work reported in this paper.

\section*{Data and Code Availability}
The complete source code, deterministic Roofline engine, 5-seed evaluation pipelines, NaN-blinded verification checks, and figure/table generation scripts supporting all findings in this study are publicly available at \url{https://github.com/PandiaJason/SHIFT-GPU}. The underlying empirical profiling dataset builds upon WattGPU \citep{fadel2026wattgpu}.

\section*{Author Biographies}

\noindent \textbf{Jason Pandian} received his B.E. and M.E. degrees in Computer Science and Engineering from Anna University, Chennai, India. He is currently an Assistant Professor in the Department of Information Technology at the Nehru Institute of Technology, Coimbatore, India, and a confirmed part-time Ph.D. research scholar under the Faculty of Computer Science and Information Engineering at Anna University, Chennai, India. Officially approved by the Doctoral Committee to proceed with his research, his doctoral thesis focuses on mitigating challenges in deep-space networking, specifically optimizing routing protocols, congestion dynamics, and delay-tolerant architectures. He has presented lightning talks on astrodynamics modeling and geospatial propagation models at the Workshop on ns-3 (WNS-3 2024, Barcelona, Spain) and the International Conference on ns-3 (ICNS-3 2025, Ritsumeikan, Japan). He is the recipient of the Best Faculty of the Year (2024) award at Nehru Institute of Technology and the Best Researcher Award (2025) at the ICNGTS International Conference. His research and instructional interests include deep-space communication routing, digital twin frameworks, programming paradigms, and data structures and algorithms (DSA).

\bibliographystyle{elsarticle-num}
\bibliography{references}

\end{document}
